\documentclass{article}
\usepackage[T1]{fontenc}
\usepackage{lmodern}
\usepackage{graphicx}
\usepackage{amsmath,amssymb}
\usepackage[a4paper,margin=2cm]{geometry}
\usepackage[absolute,overlay]{textpos}
\setlength{\TPHorizModule}{1mm}
\setlength{\TPVertModule}{1mm}
\usepackage[indonesian]{babel}
\usepackage{hyperref}
\setlength{\emergencystretch}{2em}
% unit: Halaman 10

\begin{document}

% === Halaman 10 (python/native) ===

• Contoh medan: 

 - medan bilangan bulat, <Z, +, .>

 - medan bilangan riil, <R, +, .>

 - medan bilangan rasional (p/q),  <Q, +, .>

• Sebuah medan disebut berhingga (finite field) jika himpunannya memiliki 

jumlah elemen yang berhingga. 

 Jika jumlah elemen di dalam himpunan adalah n, maka notasinya Fn

 Contoh: F2 adalah medan dengan elemen 0 dan 1

• Medan berhingga sering dinamakan juga Galois Field, untuk menghormati 

Evariste Galois, seorang matematikawan Perancis (1811 – 1832)

Rinaldi Munir/II4021 Kriptografi

10

\end{document}
